\documentclass[11pt,a4paper]{article}

\usepackage[T1]{fontenc}
\usepackage[utf8]{inputenc}
\usepackage{lmodern}
\usepackage{microtype}
\usepackage{geometry}
\usepackage{amsmath}
\usepackage{booktabs}
\usepackage{array}
\usepackage{multirow}
\usepackage{longtable}
\usepackage{xcolor}
\usepackage{enumitem}
\usepackage[hidelinks]{hyperref}
\usepackage{fancyhdr}
\usepackage{titlesec}
\usepackage{parskip}
\usepackage{amssymb}
\usepackage{tcolorbox}
\usepackage{fontawesome5}
\tcbuselibrary{skins,breakable}

%% ---- Math shorthands -------------------------------------------------------
\newcommand{\mat}[1]{\mathbf{#1}}
\newcommand{\vect}[1]{\mathbf{#1}}

%% ---- Section-title helpers (safe in bookmarks via texorpdfstring) -----------
\newcommand{\dif}[1]{\texorpdfstring{\hfill#1}{}}
\newcommand{\novel}{\texorpdfstring{\quad\textcolor{dkred}{$\star$~\textbf{Novel}}}{[Novel]}}

\geometry{top=2.5cm,bottom=2.5cm,left=2.5cm,right=2.5cm}

%% ---- Colours ----------------------------------------------------------------
\definecolor{dkblue}{RGB}{0,70,130}
\definecolor{dkgreen}{RGB}{0,110,50}
\definecolor{dkred}{RGB}{160,0,0}
\definecolor{dkorange}{RGB}{180,80,0}
\definecolor{lightblue}{RGB}{232,244,255}
\definecolor{lightgreen}{RGB}{232,255,238}
\definecolor{lightyellow}{RGB}{255,252,220}
\definecolor{lightred}{RGB}{255,235,235}
\definecolor{lightgray}{RGB}{248,248,248}
\definecolor{midgray}{RGB}{120,120,120}

%% ---- Section styling --------------------------------------------------------
\titleformat{\section}{\large\bfseries\color{dkblue}}{\thesection}{1em}{}[\color{dkblue}\titlerule]
\titleformat{\subsection}{\normalsize\bfseries\color{dkblue}}{\thesubsection}{1em}{}
\titleformat{\subsubsection}{\normalsize\bfseries\color{dkgreen}}{\thesubsubsection}{0.8em}{}

\pagestyle{fancy}
\fancyhf{}
\fancyhead[L]{\small\color{midgray}UWB RTLS --- Improvement Plan Rev\,2}
\fancyhead[R]{\small\color{midgray}\leftmark}
\fancyfoot[C]{\small\thepage}
\renewcommand{\headrulewidth}{0.4pt}

%% ---- Custom boxes -----------------------------------------------------------
\newtcolorbox{donebox}{
  enhanced,breakable,colback=lightgreen,colframe=dkgreen,
  boxrule=0.5pt,arc=3pt,left=6pt,right=6pt
}
\newtcolorbox{phasebox}[1][Phase]{
  enhanced,breakable,colback=lightblue,colframe=dkblue,
  boxrule=0.5pt,arc=3pt,left=6pt,right=6pt,
  title=\textbf{#1},fonttitle=\small\color{dkblue},
  attach boxed title to top left={yshift=-2mm,xshift=4mm},
  boxed title style={colback=white,colframe=dkblue}
}
\newtcolorbox{claimbox}{
  enhanced,breakable,colback=lightyellow,colframe=dkorange,
  boxrule=0.5pt,arc=3pt,left=6pt,right=6pt
}
\newtcolorbox{warnbox}{
  enhanced,breakable,colback=lightred,colframe=dkred,
  boxrule=0.5pt,arc=3pt,left=6pt,right=6pt
}
\newtcolorbox{codehwbox}{
  enhanced,breakable,colback=lightyellow,colframe=dkorange,
  boxrule=0.5pt,arc=3pt,left=6pt,right=6pt
}

%% ---- Difficulty badge -------------------------------------------------------
\newcommand{\easy}{\textcolor{dkgreen}{\textbf{Easy}}}
\newcommand{\medium}{\textcolor{dkorange}{\textbf{Medium}}}
\newcommand{\medhard}{\textcolor{dkorange}{\textbf{Med-Hard}}}
\newcommand{\hard}{\textcolor{dkred}{\textbf{Hard}}}
\newcommand{\vhard}{\textcolor{dkred}{\textbf{Very Hard}}}
\newcommand{\done}{\texorpdfstring{\textcolor{dkgreen}{$\checkmark$~\textbf{Done}}}{[Done]}}
\newcommand{\codehwdone}{\texorpdfstring{\textcolor{dkorange}{$\checkmark$~\textbf{Code Done (HW pending)}}}{[Code Done, HW pending]}}
\newcommand{\pendingopt}{\texorpdfstring{\textcolor{dkred}{\faExclamationTriangle~\textbf{PENDING}}}{[PENDING]}}

%% =============================================================================
\begin{document}
%% =============================================================================

\begin{titlepage}
  \centering\vspace*{2cm}
  {\color{dkblue}\rule{\linewidth}{2pt}}\\[0.5cm]
  {\huge\bfseries UWB RTLS --- Full Improvement Plan\\[0.3cm]
   \large Revision 2 \quad|\quad June 2026}\\[0.4cm]
  {\color{dkblue}\rule{\linewidth}{2pt}}\\[1.5cm]
  \begin{tabular}{ll}
    \textbf{Hardware}      & 4$\times$ Makerfabs ESP32 UWB Pro with Display (DW1000) \\
    \textbf{Configuration} & 3 anchors + 1 tag (extendable) \\
    \textbf{Goal}          & XYZ + RPY indoor positioning, no cameras, microcontroller-level \\
    \textbf{Current state} & 3--10\,cm XY accuracy, 2D only, no NLOS detection, no IMU \\
    \textbf{Target}        & $\leq$3\,cm XY, $\leq$5\,cm Z, $\leq$5$^\circ$ RPY \\
  \end{tabular}
  \vfill
\end{titlepage}

\tableofcontents
\newpage

%% =============================================================================
\section{Context and System Summary}
%% =============================================================================

The system uses Ultra-Wideband (UWB) radio time-of-flight ranging to locate a mobile tag relative to fixed anchors. The DW1000 chip achieves 15.65\,ps timestamp resolution (4.7\,mm per tick), giving a theoretical 1--2\,cm floor in LOS conditions. Current performance of 3--10\,cm is dominated by calibration errors, uncorrected multipath bias, and the absence of NLOS detection. This plan describes how to close that gap in six phases.

All position computation runs on the host PC (MATLAB). The tag firmware streams raw distances; no algorithm changes are required on the microcontroller to add more anchors or change the estimator.

%% =============================================================================
\section{Existing Library Comparison}
%% =============================================================================

\subsection{Capability Matrix}

\begin{table}[ht]
\centering
\caption{Open-source DW1000 library capability comparison (verified against current source, June 2026)}
\renewcommand{\arraystretch}{1.2}
\begin{tabular}{lcccccc}
\toprule
\textbf{Capability} & \textbf{thotro} & \textbf{Makerfabs} & \textbf{jremington} & \textbf{pizzo00} & \textbf{dw1000-ng} & \textbf{UwbRtls (ours)} \\
\midrule
Anchor ceiling       & 4         & 4         & 4$^\dagger$ & 6         & No limit$^\ddagger$ & \textbf{No limit} \\
RF collisions        & Yes       & Yes       & Yes         & Yes       & No            & \textbf{No} \\
Update rate          & 0.6\,Hz   & 0.6\,Hz   & 0.6\,Hz     & 1\,Hz     & 2\,Hz          & \textbf{3+\,Hz} \\
Multi-tag            & No        & No        & No          & Partial   & No             & \textbf{Designed-in} \\
NLOS detection       & No        & No        & No          & No        & No             & \textbf{Planned (2.1)} \\
IMU fusion           & No        & No        & No          & No        & No             & \textbf{Planned (4.4)} \\
3D + RPY             & No        & No        & Fragile     & No        & No             & \textbf{Planned (4.1--4.3)} \\
Anchor self-survey   & No        & No        & No          & No        & No             & \textbf{Planned (1.5)} \\
Actively maintained  & No (2020) & Partial   & Yes         & Yes       & No (arch.\ 2023) & \textbf{Yes} \\
\bottomrule
\end{tabular}
\end{table}

\begin{footnotesize}
$^\dagger$ jremington raises \texttt{MAX\_DEVICES} to 7 but the RANGE frame encoding still overflows at 5 anchors --- see \S\ref{sec:libdetail}.\\
$^\ddagger$ dw1000-ng uses individual unicast TWR (no shared frame), so no frame-size ceiling. Archived November 2023, single-tag only.
\end{footnotesize}

\subsection{Detailed Library Analysis}
\label{sec:libdetail}

\subsubsection{thotro/arduino-dw1000 (canonical, abandoned May 2020)}

\begin{itemize}[leftmargin=2em]
  \item \textbf{4-anchor hard ceiling.} The RANGE frame packs 17 bytes per device (2-byte address + three 5-byte absolute timestamps) into a single 90-byte buffer (\texttt{LEN\_DATA = 90}). With the 9-byte MAC header: $\lfloor(90-9-2)/17\rfloor = 4$ devices. Raising \texttt{MAX\_DEVICES} past 4 causes a silent buffer overrun --- confirmed in Issue \#81. No bounds checking.
  \item \textbf{Broadcast POLL collisions.} All anchors attempt to reply to POLL\_ACK simultaneously; only one wins per exchange, starving others.
  \item \textbf{Volatile bug.} \texttt{\_networkDevicesNumber} not declared \texttt{volatile}; compiler may cache stale value, corrupting the device list (Issue \#179, unfixed).
  \item \textbf{NLOS data unused.} \texttt{getFirstPathPower()} and \texttt{getReceivePower()} are implemented and stored on each device object, but \textbf{the difference is never computed, no NLOS score is derived, and no weighting or rejection logic exists}.
\end{itemize}

\subsubsection{Makerfabs mf\_DW1000}
Thotro fork with ESP32 pin remapping. Law-of-cosines solver hardcoded for exactly 3 anchors. All thotro limitations apply unchanged.

\subsubsection{jremington/UWB-Indoor-Localization\_Arduino}

The repo bundles two distinct library copies:

\textbf{DW1000\_library} (thotro derivative): Adds \texttt{volatile} fix to \texttt{\_networkDevicesNumber}; raises \texttt{MAX\_DEVICES} to 7. However, the RANGE frame encoding is \textbf{identical} --- still 17 bytes/device, $\lfloor79/17\rfloor = 4$ fit before overflow. Setting \texttt{MAX\_DEVICES = 7} without changing the frame is a latent buffer overrun: devices 5--7 silently corrupt memory.

\textbf{DW1000\_library\_pizzo00} (Pizzolato fork, hosted here since pizzo00's original repo is deleted): Reduces per-device bytes to 12 using \textbf{differential} timestamps (\texttt{rangeDeviceSize = 12}), capping at \texttt{devicePerTransmit = 6} per frame ($\lfloor79/12\rfloor = 6$). This fixes the frame-size arithmetic for DW1000. Still uses broadcast POLL (RF collisions remain). Includes on-device LLS solver with precomputed $(A^\top A)^{-1}$ (good technique, borrowed in our Phase 6.3).

\subsubsection{F-Army/arduino-dw1000-ng (archived November 2023)}

\textbf{Architecture is fundamentally different}: individual unicast TWR per anchor (no shared broadcast frame). The \texttt{tagRangeInfrastructure()} function sequences anchors one at a time via a \texttt{next\_anchor} field in each anchor's confirm reply. Consequently, there is no frame-size anchor ceiling. However: single-tag only, no NLOS detection, archived and unmaintained.

\begin{claimbox}
\textbf{Corrected novelty claim on anchor scaling.} The thotro/Makerfabs/jremington lineage has a \textbf{hard 4-anchor overflow boundary} in the shared RANGE frame. The pizzo00 fork fixes the frame math to reach 6. The F-Army dw1000-ng library also has no anchor ceiling via individual unicast TWR, but is archived (2023), single-tag only, and has no NLOS detection or IMU fusion.

\textbf{Our system is the first to combine:} (1) individual addressed DS-TWR with no anchor ceiling, (2) multi-tag TDMA design, (3) FP\_POWER NLOS detection, and (4) IMU fusion --- in a single open-source library.
\end{claimbox}

\subsubsection{NLOS Detection --- Verified Novelty}

\begin{claimbox}
\textbf{Confirmed: no existing open-source DW1000 positioning library uses FP\_POWER for NLOS detection.}

Every audited library (thotro, Makerfabs, jremington/DW1000\_library, jremington/DW1000\_library\_pizzo00, F-Army/dw1000-ng, linorobot/ros\_dwm1000) calls \texttt{getFirstPathPower()} and \texttt{getReceivePower()} and stores both values on the device object --- but \textbf{none computes their difference, none derives an NLOS score, and none uses that score to weight or reject a measurement.}

The F-Army library has an \texttt{\_nlos} boolean flag, but it controls the LDE leading-edge algorithm configuration (a static chip setting), not per-packet NLOS inference.

Academic papers (e.g.\ arXiv:2403.19706) use the APS006 FP\_power metric in post-processing Python/MATLAB scripts, not in the embedded Arduino positioning library.

\textbf{Our implementation of per-packet} $\Delta P_\text{NLOS} = P_\text{RX} - P_\text{FP}$ \textbf{with weighted multilateration is, to our knowledge, the first in an open-source DW1000 positioning library.}
\end{claimbox}

\subsection{Techniques Borrowed from Existing Libraries}

\begin{table}[ht]
\centering
\caption{Techniques incorporated from existing work}
\begin{tabular}{lll}
\toprule
\textbf{Source} & \textbf{Technique} & \textbf{Our phase} \\
\midrule
jremington & Age-gated anchor eviction (MATLAB) & Phase 0 (\done) \\
jremington / pizzo00 & On-device LLS with precomputed $(A^\top A)^{-1}$ & Phase 6.3 \\
pizzo00 & Differential timestamp encoding (12\,B vs 17\,B/device) & Phase 6.2 \\
dw1000-ng & Deep sleep on anchors between ranging cycles & Phase 5.1 \\
Nobody & FP\_POWER / RX\_POWER NLOS score \& weighting & Phase 2.1 (\textbf{novel}) \\
Nobody & Anchor self-survey via MDS & Phase 1.5 (\textbf{novel}) \\
Nobody & Full 6-DOF (XYZ + RPY) with UWB + IMU fusion & Phase 4 (\textbf{novel}) \\
\bottomrule
\end{tabular}
\end{table}

%% =============================================================================
\section{Implementation Status}
%% =============================================================================

\begin{donebox}
\textbf{Phase 0 --- Already Complete}
\begin{itemize}[leftmargin=2em,topsep=4pt,itemsep=2pt]
  \item[$\checkmark$] MATLAB watchdog + UDP auto-reconnect (10\,s timeout $\to$ recreates \texttt{udpport})
  \item[$\checkmark$] EKF divergence guard + auto-reset (jump $>$2\,m or $\mathrm{tr}(P) > 25$ triggers reinit)
  \item[$\checkmark$] Age-gated anchor eviction in MATLAB (\texttt{AGE\_GATE\_SEC = 3.0\,s})
  \item[$\checkmark$] \texttt{RANGE\_BIAS\_M} tuning slider (0--0.10\,m global multipath bias subtraction)
  \item[$\checkmark$] Inner try-catch (bad sweeps skip silently, session survives)
  \item[$\checkmark$] Status label in tuning panel (packet count, EKF resets, current bias)
  \item[$\checkmark$] Live 5-slider tuning panel with Stationary / Slow Walk / Fast Walk presets
\end{itemize}
\end{donebox}

%% =============================================================================
\section{Phase 1 --- Foundation Fixes}
%% =============================================================================

\begin{phasebox}[Phase 1 --- Foundation | No algorithm changes. Highest ROI for effort.]
Attacking calibration, protocol, and configuration issues before adding complexity.
Expected improvement: 3--10\,cm $\to$ 2--6\,cm XY.
\end{phasebox}

\subsection{1.0 Switch to 64\,MHz PRF Accuracy Mode\dif{\easy}}

\begin{codehwbox}
\textbf{Status (2026-06-10):} \codehwdone{} --- \texttt{UWB\_RADIO\_MODE} changed to \texttt{MODE\_LONGDATA\_RANGE\_ACCURACY} and \texttt{UWB\_REPLY\_DELAY\_US} reduced 7000\,$\to$\,6000\,µs in \texttt{UwbConfig.h} (single-source, applies to all boards). \textbf{Hardware follow-up:} reflash all 4 boards, re-run calibration (antenna delay values change between PRF modes).
\end{codehwbox}

\textbf{Change:} \texttt{DW1000.enableMode(MODE\_LONGDATA\_RANGE\_LOWPOWER)} $\to$ \texttt{MODE\_LONGDATA\_RANGE\_ACCURACY} in \texttt{UwbConfig.h}.

\textbf{Why:} 64\,MHz PRF produces a sharper autocorrelation peak in the channel impulse response. The receiver can more reliably identify the first-path arrival and reject multipath peaks arriving slightly later. The Pro boards already have a PA/LNA amplifier for range --- this switch improves accuracy, not range.

\textbf{Actions:} Reflash all 4 boards. Reduce reply delay \texttt{UWB\_REPLY\_DELAY\_US} 7000 $\to$ 6000\,µs. Re-run calibration (1.1) --- delay register values differ between PRF modes.

\textbf{Removes:} 1--3\,cm systematic ranging error from multipath peak ambiguity.

\subsection{1.1 Better Calibration\dif{\easy}}

\begin{codehwbox}
\textbf{Status (2026-06-10):} \codehwdone{} --- \texttt{SAMPLES\_PER\_STEP} 40\,$\to$\,200, \texttt{SEARCH\_ITERS} 12\,$\to$\,14, and loop counter widened to \texttt{uint16\_t} (prevents silent wrap at $>$255) in \textbf{both} \texttt{AntennaCalibration.ino} copies (\texttt{sketches/} active + \texttt{libraries/.../examples/} reference). Post-calibration loop throttled to 1\,Hz; calibrated delay printed on every line. \textbf{Hardware follow-up:} run calibration on each board, validate at two distances ($\approx$1.5\,m and $\approx$5\,m, agree within $\pm$2\,cm).
\end{codehwbox}

\textbf{Change:} \texttt{SAMPLES\_PER\_STEP} 40 $\to$ 200; \texttt{SEARCH\_ITERS} 12 $\to$ 14 in \texttt{AntennaCalibration.ino}.

\textbf{Why:} 40 samples gives $\pm$1\,cm statistical uncertainty on the mean, propagating to $\pm$3 delay-register units of systematic error. 200 samples reduces this to $\pm$0.3\,cm.

\textbf{Validation:} After convergence run the calibration loop for 60\,s and record mean $\pm$ std. Validate at two distances (e.g.\ 1.5\,m and 5\,m) --- both should agree within $\pm$2\,cm.

\textbf{Removes:} $\pm$1--3\,cm systematic calibration bias per board.

\subsection{1.2 Adaptive Polling + Age-Gated Anchor Eviction\dif{\easy}}

\begin{donebox}
\textbf{Status (2026-06-11):} \done{} --- \texttt{UwbScheduler} now tracks per-anchor \texttt{\_failStreak}, \texttt{\_skipSweeps}, and \texttt{\_backoffMult}. After 3 consecutive failures the anchor is skipped; skip duration doubles each subsequent timeout: 5\,$\to$\,10\,$\to$\,20\,$\to$\,40 sweeps (capped). Serial prints \texttt{[SCHED]} on skip and on recovery. MATLAB age-gated eviction was already complete. Library-only change --- no reflash required.
\end{donebox}

\textbf{Firmware (remaining):} Add \texttt{failCount[N]} per anchor in \texttt{UwbScheduler}. After 3 consecutive timeouts, skip that anchor for 5 sweeps then retry with exponential backoff. This avoids wasting 100\,ms per dead anchor per sweep.

\textbf{MATLAB:} \done{} --- \texttt{anchorLastSeen} struct with \texttt{tic}/\texttt{toc} tracking; \texttt{AGE\_GATE\_SEC = 3.0} evicts stale range history automatically.

\subsection{1.3 Faster Sweep Rate\dif{\medium}}

\begin{codehwbox}
\textbf{Status Option A (2026-06-11):} \codehwdone{} --- \texttt{UWB\_REPLY\_DELAY\_US} reduced 6000\,$\to$\,5000\,µs (previous was 7000\,µs before Phase\,1.0). Estimated sweep rate improvement: $\sim$8\,Hz\,$\to$\,$\sim$10\,Hz (4 anchors). Minimum safe value for 110\,kbps / 64\,MHz PRF mode is $\approx$3500\,µs; 5000\,µs leaves $\approx$1.5\,ms margin above ESP32 SPI + DW1000 processing time. \textbf{Hardware follow-up:} reflash all boards; if exchange failure rate increases, revert to 6000\,µs.
\end{codehwbox}

\begin{warnbox}
\textbf{Option B --- \pendingopt{}} Switch to \texttt{MODE\_SHORTDATA\_FAST\_LOWPOWER} (6.8\,Mbps, 64-symbol preamble, reply delay $\sim$1500\,µs). Estimated $\sim$35--45\,Hz for 4 anchors. \textbf{Deferred:} trades 64\,MHz PRF accuracy for speed --- increases multipath sensitivity in small rooms. Defer until Phase\,2.1 NLOS detection is in place to compensate.
\end{warnbox}

\textbf{Option A (safe):} Reduce \texttt{UWB\_REPLY\_DELAY\_US} 6000 $\to$ 5000\,µs. Saves 2\,ms per exchange, $\sim$8\,ms per sweep with 4 anchors.

\textbf{Option B (aggressive):} Switch to \texttt{MODE\_SHORTDATA\_FAST\_LOWPOWER} (6.8\,Mbps, 64-symbol preamble). Reply delay drops to $\sim$1500\,µs giving $\sim$35--45\,Hz for 4 anchors. Cost: increased multipath sensitivity (16\,MHz PRF vs 64\,MHz PRF).

Higher sweep rate directly improves all filter layers downstream.

\subsection{1.4 Multipath Bias Correction\dif{\done}}

Global \texttt{RANGE\_BIAS\_M} subtracted from all filtered ranges before the multilaterator. Tunable 0--0.10\,m via slider. Tune by placing tag at a known point and minimising \texttt{info.rms}.

\subsection{1.5 Anchor Self-Survey (Auto-Localization)\dif{\medium}}

\textbf{What:} Anchors range to each other. MATLAB computes the anchor layout via Multidimensional Scaling (MDS) and writes \texttt{anchors.json} automatically. No tape measure required.

\textbf{Firmware:} New sketch \texttt{examples/AnchorSurvey/AnchorSurvey.ino}. Each anchor temporarily takes the tag role and ranges to all other anchor addresses. Sends \texttt{SURVEY,v1,<src>,<dst>,<dist\_mm>,<quality>} over HostLink.

\textbf{MATLAB:} New script \texttt{matlab/runSurvey.m}. Collects $\binom{N}{2}$ pairwise distances (averaged over 100 sweeps per pair). Runs classical MDS: build squared-distance matrix $D^{(2)}$, apply double-centering $B = -\tfrac{1}{2}JD^{(2)}J$, eigendecompose, take top-$d$ eigenpairs. Fix frame (anchor 1 at origin, anchor 2 on $+X$, anchor 3 in $+Y$ half-plane). Write \texttt{anchors.json}.

\textbf{Accuracy:} $\approx$1--2\,cm anchor position accuracy after 100-sample averaging ($\sigma_d/\sqrt{100} \approx 0.5$\,cm per pair).

\textbf{Limitation:} 3D self-survey requires meaningful height variation. All-same-height gives a degenerate $Z$ axis.

%% =============================================================================
\section{Phase 2 --- Ranging Quality}
%% =============================================================================

\begin{phasebox}[Phase 2 --- Ranging Quality | Firmware + MATLAB. Attacks measurement errors at source.]
Expected improvement: 2--6\,cm $\to$ 1--4\,cm XY.
\end{phasebox}

\subsection{2.1 FP\_POWER NLOS Detection\dif{\medium}\novel}

\textbf{What:} Read first-path amplitude registers after each DW1000 receive. Compute per-packet NLOS score per Decawave APS006. Send in RANGE\_REPORT payload. Parse and act in MATLAB.

\textbf{Registers:} FP\_AMPL1, FP\_AMPL2, FP\_AMPL3 (reg 0x12, offsets 0x07--0x0C); RXPACC (reg 0x10, offset 0x2C); CIR\_PWR (reg 0x12, offset 0x00).

\textbf{Math:}
\begin{align*}
P_\text{FP} &= 10\log_{10}\!\left(\frac{A_1^2+A_2^2+A_3^2}{N_\text{acc}^2}\right) - A_\text{const}
\quad (A_\text{const} = 113.77\;\text{dBm at 64\,MHz PRF}) \\
\Delta P_\text{NLOS} &= P_\text{RX} - P_\text{FP}
\end{align*}

\begin{center}
\begin{tabular}{cll}
\toprule
\textbf{Score} & \textbf{Condition} & \textbf{Action} \\
\midrule
$<3$\,dB & LOS & Full weight \\
3--6\,dB & Soft NLOS & Reduce weight \\
$>6$\,dB & Hard NLOS & Strong downweight or reject \\
\bottomrule
\end{tabular}
\end{center}

\textbf{Why it matters:} Detects NLOS before the solver runs, based on physics. Works even when consistent NLOS bias looks ``normal'' to MAD gating.

\textbf{Novelty:} Confirmed by auditing all major DW1000 libraries (thotro, Makerfabs, jremington, pizzo00, dw1000-ng, linorobot). Every library exposes \texttt{getFirstPathPower()} and \texttt{getReceivePower()} but \textbf{none computes their difference, derives an NLOS score, or uses it to weight or reject measurements.}

\subsection{2.2 Weighted Multilateration\dif{\medium}}

Derive weights from NLOS score:
\[
\sigma_i = \sigma_\text{base} \cdot 10^{\Delta P_{\text{NLOS},i}/20}, \qquad w_i = 1/\sigma_i^2
\]
A 10\,dB NLOS anchor receives $1/100$ the weight of a clean LOS anchor. Modify Levenberg-Marquardt solver in \texttt{Multilaterator.m} to use weighted residuals and Jacobian $\mat{J}^\top\mat{W}\mat{J}$.

\subsection{2.3 NLOS Range Bias Correction\dif{\easy}}

Physical bias correction applied per anchor before the solver:
\[
d_\text{corrected} = d_\text{measured} - k_\text{NLOS} \cdot \max(0,\; \Delta P_\text{NLOS} - 3\;\text{dB})
\]
Tunable slider \texttt{K\_NLOS} (0--0.06\,m/dB) in the tuning panel. Removes systematic positive through-wall ranging bias (typically 20--50\,cm per obstructed wall).

\textbf{Requires:} Phase 2.1 for the per-packet NLOS score.

\subsection{2.4 Per-Anchor Diagnostics Panel\dif{\easy}}

Second axes in \texttt{LivePlotter} showing per-anchor range residual (bar chart) and NLOS score (colour-coded: green/yellow/red) updated every sweep. Provides immediate visual feedback on which anchor is causing problems.

\subsection{\texorpdfstring{2.5 GPS-Style $\delta$-Solve (4+ Anchors)\dif{\medium}}{2.5 GPS-Style delta-Solve (4+ Anchors)}}

Augment the multilateration state to $[\vect{x}^\top, \delta]^\top$ where $\delta$ is a common range bias:
\[
\hat{d}_i = \|\vect{x} - \vect{a}_i\| + \delta + e_i
\]
Jacobian row $i$: $[\text{unit\_vec}_i,\; 1]$. Solved $\delta$ is a real-time calibration health metric.

\begin{warnbox}
\textbf{Warning:} With $N = d+1$ anchors (3 in 2D) this reduces to TDOA geometry --- ill-conditioned near the anchor centroid. Only enable with $N \geq d+2$ anchors.
\end{warnbox}

\textbf{Requires:} Phase 4.1 (4th anchor).

%% =============================================================================
\section{Phase 3 --- State Estimation}
%% =============================================================================

\begin{phasebox}[Phase 3 --- State Estimation | Better filter models. MATLAB only.]
Expected improvement: 1--4\,cm $\to$ 1--3\,cm XY.
\end{phasebox}

\subsection{\texorpdfstring{3.1 Constant-Acceleration EKF (CV $\to$ CA Model)\dif{\medium}}{3.1 Constant-Acceleration EKF (CV to CA Model)}}

Upgrade from 4-state $[\vect{p}, \vect{v}]$ to 6-state $[\vect{p}, \vect{v}, \vect{a}]$. Dynamics: position $\leftarrow$ position $+$ velocity $\cdot\Delta t + \tfrac{1}{2}$ acceleration $\cdot\Delta t^2$; acceleration $\leftarrow$ acceleration $+$ random jerk (Singer model). Better tracks direction changes and acceleration phases; reduces lag on turns.

\subsection{\texorpdfstring{3.2 $\delta$ as EKF State (Works with 3 Anchors)\dif{\medhard}}{3.2 delta as EKF State (Works with 3 Anchors)}}

Add $\delta$ to EKF state: $[\vect{p}, \vect{v}, \delta]$ or $[\vect{p}, \vect{v}, \vect{a}, \delta]$. $\delta$ evolves as a random walk. Tight coupling: raw ranges feed directly into EKF measurement update. Measurement: $h_i(\vect{x}) = \|\vect{p}-\vect{a}_i\| + \delta$. Jacobian row $i$: $[\hat{u}_i,\; \mathbf{0},\; 1]$ (unit vector + zeros for velocity/accel + 1 for bias).

\textbf{Advantage over multilaterator $\delta$-solve:} Works with 3 anchors; temporally smoothed; avoids ill-conditioning near centroid.

\subsection{3.3 Moving Horizon Estimator (MHE)\dif{\hard}}

Constrained optimisation over a sliding window of $N$ past sweeps. At time $k$, minimise:
\[
\min_{\{\vect{x}_t\}}\; \underbrace{\|\vect{x}_{k-N+1} - \bar{\vect{x}}_{k-N+1}\|^2_{\mat{P}_\text{arr}^{-1}}}_{\text{arrival cost}}
+ \underbrace{\sum_t \|\vect{x}_t - \mat{F}\vect{x}_{t-1}\|^2_{\mat{Q}^{-1}}}_{\text{dynamics}}
+ \underbrace{\sum_t\sum_i \rho_\delta\!\left(\|\vect{p}_t-\vect{a}_i\| - \hat{d}_{t,i}\right)}_{\text{range residuals (Huber)}}
\]
subject to $\vect{p}_\text{min} \leq \vect{p}_t \leq \vect{p}_\text{max}$ (room boundary) and $\|\vect{v}_t\| \leq v_\text{max}$ (speed limit).

Huber loss threshold $\delta = 0.15$\,m. Computation: \texttt{fmincon} with $N=6$, $d=4$ states $\Rightarrow$ 2--5\,ms per solve.

\textbf{Advantages over EKF:} NLOS absorbed by Huber loss; room boundary prevents wild jumps; physically impossible positions excluded.

\subsection{3.4 Hybrid EKF + MHE\dif{\medium}}

EKF every sweep (real-time, zero lag). MHE every $K=5$ sweeps (accuracy pass). If outputs diverge beyond threshold, flag increased uncertainty. Requires Phase 3.3 complete.

%% =============================================================================
\section{Phase 4 --- 3D + Orientation (XYZ + RPY)}
%% =============================================================================

\begin{phasebox}[Phase 4 --- 3D + RPY | Firmware + MATLAB. Adds the 3rd dimension and full orientation.]
Expected outcome: XY 1--3\,cm, Z 2--5\,cm, RPY 3--5$^\circ$ at 100\,Hz.
\end{phasebox}

\subsection{4.1 4th Anchor + 3D Configuration\dif{\easy}}

Flash board 4 as \texttt{ANCHOR\_ID = 0x04}. Mount at a different height from the others. Update \texttt{anchors.json}: add entry with non-zero $z$, set \texttt{"dim": 3}. Add \texttt{0x04} to \texttt{ANCHORS[]} in \texttt{Tag.ino}. This also enables the $\delta$-solve (Phase 2.5) since we now have $d+2 = 5$ constraints.

\subsection{4.2 BNO085 IMU --- Firmware Implementation\dif{\medium}}

Implement \texttt{SensorImu::begin()} and \texttt{SensorImu::read()} using the Adafruit BNO08x SPI library.

\textbf{Why SPI:} The I2C bus (pins 4/5) is occupied by the SSD1306 OLED. The BNO085 also has known clock-stretching issues on ESP32 I2C.

\textbf{Reports:} ARVR Stabilised Rotation Vector at 10\,ms intervals (quaternion, 100\,Hz) + Linear Acceleration at 10\,ms (gravity removed, 100\,Hz). The IMU tail in HostLink and the parser in \texttt{FrameParser.m} are already wired for this data.

\subsection{4.3 RPY Output from Quaternion\dif{\easy}}

Extract Euler angles from the BNO085 quaternion $[q_w, q_x, q_y, q_z]$ using the ZYX convention:
\begin{align*}
\phi   &= \text{atan2}\!\left(2(q_wq_x + q_yq_z),\; 1-2(q_x^2+q_y^2)\right) && \text{(roll)} \\
\theta &= \arcsin\!\left(2(q_wq_y - q_zq_x)\right) && \text{(pitch)} \\
\psi   &= \text{atan2}\!\left(2(q_wq_z + q_xq_y),\; 1-2(q_y^2+q_z^2)\right) && \text{(yaw)}
\end{align*}
Log RPY alongside XYZ. Add RPY readout to HUD. The magnetometer in the BNO085 provides an absolute yaw reference, preventing unbounded yaw drift.

\textbf{Requires:} Phase 4.2.

\subsection{4.4 Loose Coupling --- IMU-Aided EKF\dif{\hard}}

IMU predict step at 100\,Hz (fills the 333\,ms gaps between UWB sweeps). UWB measurement update at 3\,Hz corrects accumulated drift.

\textbf{State:} $[\vect{p}, \vect{v}, q_w, q_x, q_y, q_z, \delta]$ (11-state, 3D). IMU prediction:
\begin{align*}
\vect{a}^w &= \mat{R}(\vect{q})\,\vect{a}^b \quad\text{(rotate to world frame)} \\
\vect{p}_{k+1} &= \vect{p}_k + \vect{v}_k\Delta t + \tfrac{1}{2}\vect{a}^w_k\Delta t^2 \\
\vect{q}_{k+1} &= \vect{q}_k \otimes \exp\!\left(\tfrac{\Delta t}{2}\vect{\omega}^b_k\right)
\end{align*}

Effectively decouples position accuracy (UWB) from position smoothness (IMU).

\subsection{4.5 Tight Coupling --- Raw Ranges + IMU in One Filter\dif{\vhard}}

\textbf{State:} $[\vect{p}, \vect{v}, \vect{q}, \delta]$ (11-state). IMU drives prediction at 100\,Hz. Raw UWB ranges (not multilaterated position) are the measurement at 3\,Hz. Solver: MHE with Huber loss, room bounds, speed limit, and quaternion unit-norm constraint $\|\vect{q}\| = 1$.

This is the architecture commercial UWB positioning systems use on dedicated hardware. Full 6-DOF output: XYZ + RPY + velocity + bias estimate.

%% =============================================================================
\section{Phase 5 --- System Hardening}
%% =============================================================================

\subsection{5.1 Deep Sleep on Anchors (from dw1000-ng)\dif{\medium}}
Put DW1000 into deep sleep for $\sim$80\% of the sweep period between ranging cycles. Reduces anchor current from $\sim$100\,mA continuous to $\sim$20\,mA average --- significant for battery-powered deployment.

\subsection{5.2 DW1000 Driver Error Handling\dif{\medium}}
The vendored driver contains 40+ \texttt{// TODO proper error/warning handling} comments. Add recovery for SPI lockup, PLL failure, and LDE timeout. Auto-reset radio if it goes silent for $>$500\,ms.

\subsection{5.3 Temperature Compensation\dif{\medium}}
The DW1000 on-chip temperature sensor is OTP-calibrated and already read by the driver but never applied. Clock drift is approximately 0.1 ticks/$^\circ$C $\approx$ 0.5\,mm/$^\circ$C per board. Apply a linear correction relative to the calibration temperature.

\subsection{5.4 Online Self-Calibration\dif{\hard}}
Batch optimisation jointly solving for anchor positions and per-board delays using accumulated tag trajectory + range data. Requires spatial diversity (tag must visit different areas of the room). Replaces the manual calibration procedure entirely.

%% =============================================================================
\section{Phase 6 --- Scale and Future}
%% =============================================================================

\subsection{6.1 Multi-Tag TDMA Superframe\dif{\hard}}
Time-slotted TDMA. Master anchor broadcasts slot assignments via ANNOUNCE frames (type byte already defined in \texttt{UwbFrame.h}). MATLAB maintains a separate EKF/MHE per \texttt{tagId}. Requires careful timing budget: $N_\text{anchors} \times N_\text{tags}$ TWR exchanges per superframe.

\subsection{6.2 Differential Timestamp Encoding (pizzo00 technique)\dif{\hard}}
If a broadcast POLL mode is ever needed (e.g.\ for firmware-side solving), adopt the pizzo00 12-byte/device encoding (two differential timestamps instead of three absolute) to reach 6 anchors per RANGE frame.

\subsection{6.3 On-Device 2D Solve (jremington technique)\dif{\medium}}
Precompute $(A^\top A)^{-1}$ once at startup (anchors are fixed). Position update becomes a matrix-vector multiply. Display $(x, y)$ on the OLED without a host PC. Useful for standalone deployment.

%% =============================================================================
\section{Full Item Index}
%% =============================================================================

\begin{longtable}{llll}
\toprule
\textbf{Item} & \textbf{Description} & \textbf{Difficulty} & \textbf{Status} \\
\midrule
\endhead
\multicolumn{4}{l}{\textit{Phase 0 --- Already Done}} \\
\midrule
0.1 & MATLAB watchdog + UDP auto-reconnect & --- & \done \\
0.2 & EKF divergence guard + auto-reset & --- & \done \\
0.3 & Age-gated anchor eviction (MATLAB) & --- & \done \\
0.4 & \texttt{RANGE\_BIAS\_M} tuning slider & --- & \done \\
0.5 & Inner try-catch (bad sweeps skip, session survives) & --- & \done \\
0.6 & Status label in tuning panel & --- & \done \\
\midrule
\multicolumn{4}{l}{\textit{Phase 1 --- Foundation}} \\
\midrule
1.0 & Switch to 64\,MHz PRF (\texttt{MODE\_LONGDATA\_RANGE\_ACCURACY}) & \easy & \codehwdone \\
1.1 & Better calibration (200 samples, multi-distance) & \easy & \codehwdone \\
1.2 & Adaptive polling --- skip dead anchors (firmware) & \easy & \done \\
1.3 & Faster sweep rate --- Option A \codehwdone; Option B \pendingopt & \medium & Partial \\
1.5 & Anchor self-survey + auto \texttt{anchors.json} & \medium & Pending \\
\midrule
\multicolumn{4}{l}{\textit{Phase 2 --- Ranging Quality}} \\
\midrule
2.1 & FP\_POWER NLOS detection (firmware + wire format) & \medium & Pending \\
2.2 & Weighted multilateration (NLOS score $\to$ weights) & \medium & Pending \\
2.3 & NLOS range bias correction slider & \easy & Pending \\
2.4 & Per-anchor diagnostics panel in LivePlotter & \easy & Pending \\
2.5 & GPS-style $\delta$-solve in Multilaterator (4+ anchors) & \medium & Pending \\
\midrule
\multicolumn{4}{l}{\textit{Phase 3 --- State Estimation}} \\
\midrule
3.1 & Constant-acceleration EKF (CV $\to$ CA model) & \medium & Pending \\
3.2 & $\delta$ as EKF state (bias estimation, 3 anchors OK) & \medhard & Pending \\
3.3 & MHE (sliding window, Huber loss, room constraints) & \hard & Pending \\
3.4 & Hybrid EKF + MHE & \medium & Pending \\
\midrule
\multicolumn{4}{l}{\textit{Phase 4 --- 3D + RPY}} \\
\midrule
4.1 & 4th anchor + 3D \texttt{anchors.json} & \easy & Pending \\
4.2 & BNO085 \texttt{SensorImu} implementation (SPI) & \medium & Pending \\
4.3 & RPY output from quaternion & \easy & Pending \\
4.4 & Loose coupling --- IMU-aided EKF (100\,Hz predict) & \hard & Pending \\
4.5 & Tight coupling --- raw ranges + IMU in MHE & \vhard & Pending \\
\midrule
\multicolumn{4}{l}{\textit{Phase 5 --- System Hardening}} \\
\midrule
5.1 & Deep sleep on anchors & \medium & Pending \\
5.2 & DW1000 driver error handling & \medium & Pending \\
5.3 & Temperature compensation & \medium & Pending \\
5.4 & Online self-calibration & \hard & Pending \\
\midrule
\multicolumn{4}{l}{\textit{Phase 6 --- Scale}} \\
\midrule
6.1 & Multi-tag TDMA superframe & \hard & Pending \\
6.2 & Differential timestamp encoding & \hard & Pending \\
6.3 & On-device 2D solve (OLED display) & \medium & Pending \\
\bottomrule
\end{longtable}

%% =============================================================================
\section{Expected Accuracy at Each Phase}
%% =============================================================================

\begin{table}[ht]
\centering
\caption{Expected accuracy progression through implementation phases}
\renewcommand{\arraystretch}{1.3}
\begin{tabular}{cllllp{5cm}}
\toprule
\textbf{After} & \textbf{XY} & \textbf{Z} & \textbf{RPY} & \textbf{Rate} & \textbf{Key enabler} \\
\midrule
Current  & 3--10\,cm & ---      & ---         & 3\,Hz       & Calibration noise + NLOS + CV EKF \\
Phase 1  & 2--6\,cm  & ---      & ---         & 4--8\,Hz    & Calibration bias removed, PRF improved \\
Phase 2  & 1--4\,cm  & ---      & ---         & 4--8\,Hz    & NLOS weighted/rejected before solver \\
Phase 3  & 1--3\,cm  & ---      & ---         & 4--8\,Hz    & Better filter model, MHE room constraints \\
Phase 4  & 1--3\,cm  & 2--5\,cm & 3--5$^\circ$ & 100\,Hz RPY & IMU fills gaps, 3D anchor geometry \\
Phase 5+ & 1--3\,cm  & 2--5\,cm & 3--5$^\circ$ & 100\,Hz RPY & Reliability, power, robustness \\
\bottomrule
\end{tabular}
\end{table}

%% =============================================================================
\section{Recommended Implementation Sequence}
%% =============================================================================

\begin{center}
\renewcommand{\arraystretch}{1.4}
\begin{tabular}{ccccccccccccc}
\textcolor{dkgreen}{\textbf{1.0}} & $\to$ & \textcolor{dkgreen}{\textbf{1.1}} & $\to$ & \textcolor{dkgreen}{\textbf{1.2}} & $\to$ & \textcolor{dkgreen}{\textbf{1.3A}} & $\to$ & 1.5 & $\to$ & 2.1 & $\to$ & 2.2 \\
\textcolor{dkgreen}{PRF} && \textcolor{dkgreen}{Cal} && \textcolor{dkgreen}{Poll} && \textcolor{dkgreen}{Rate-A} && Survey && NLOS && Wgt \\[4pt]
$\to$ & 2.4 & $\to$ & 3.1 & $\to$ & 4.1 & $\to$ & 4.2 & $\to$ & 4.3 & $\to$ & 4.4 & \\
& Plot && CA-EKF && 4th Anchor && IMU && RPY && Fusion &
\end{tabular}
\end{center}

\medskip
\noindent\textcolor{dkgreen}{Green} = code complete. \textcolor{dkorange}{Orange} = code done, hardware follow-up pending.
\textcolor{dkred}{Red} items (1.3B fast mode) deferred until Phase\,2.1 NLOS detection is in place.

\bigskip
Phases 2.5, 3.2, 3.3, 3.4, 4.5, and all of Phases 5--6 can be added in any order after 4.4 is stable.

\end{document}
